
\section{第 3 章概述}
\label{sec:3-1}

正如上一章所说明的，一般而言，拉格朗日交弗洛尔上同调无法定义。困难来自伪全纯条带的边界上会冒出伪全纯圆盘。若按通常方式，用某些伪全纯圆盘或条带之模空间的稳定映射紧化来定义“上边界算子”，冒出的圆盘便会出现在实余维一的边界分支中。另一方面，任意闭辛流形上的周期哈密顿系统都有可定义的弗洛尔上同调
\textbf{[FuOn99II, LiuTi98, Rua99]}。在这种情形下，用来定义上边界算子的是无边界亏格零稳定映射的模空间；冒出的球面出现在实余维二（即复余维一）的层中。因此，一旦借助 \textbf{[FuOn99II]} 引入的仓西结构解决横截性问题，就能证明这些冒出的球面不会对上边界算子作出贡献，理论因而可以在（上）同调层面表述。拉格朗日交的情形则不同：与圆盘冒泡对应的实余维一层迫使我们在（上）链层面发展理论。

为此，我们引入并构造一种称为\emph{滤过 $A_\infty$ 代数}的同伦代数结构。粗略地说，一个 $A_\infty$ 代数是一个分次模 $\bigoplus C^\bullet$，连同无穷多个运算
$\mathfrak m=\{\mathfrak m_k\}_{k=0,1,2,\ldots}$，其中
\[
  \mathfrak m_k\colon
  \underbrace{C^*\otimes\cdots\otimes C^*}_{k\text{ 个因子}}
  \longrightarrow C^*,
  \qquad k=0,1,2,\ldots,
\]
并且这些运算满足称为 $A_\infty$ 公式或 $A_\infty$ 关系的若干关系。我们称
$\mathfrak m$ 为一个 $A_\infty$ 结构。精确定义见定义 3.2.3（无滤过
$A_\infty$ 代数）和定义 3.2.20（滤过 $A_\infty$ 代数）。第 3.2 节将介绍
$A_\infty$ 代数与滤过 $A_\infty$ 代数的基本材料。

从历史上说，斯塔谢夫受到菅原正博 \textbf{[Sug57]} 的工作启发，在
\textbf{[Sta63]} 中引入了（无滤过）$A_\infty$ 代数的概念，以研究环路空间的同伦论；相关工作另见 \textbf{[Ada78, BoVo73]}。$A_\infty$ 结构的一个几何侧面，是它描述树（即可缩的一维 CW 复形）的余维一退化。全纯圆盘（亏格为零）冒泡的所有构形都能用树来描述，所以 $A_\infty$ 结构会自然出现在这里。$A_\infty$ 代数可以看作微分分次代数（D.G.A.）的推广，并且在形变理论中起着重要作用，例如见 \textbf{[GoMi88,90, GugSt86]}。我们将用相应的 $A_\infty$ 代数研究拉格朗日子流形及其弗洛尔上同调的形变理论；这是第 4、5 章的主题。

在陈述本章主要结果之前，先说明定向问题。带拉格朗日边界条件的伪全纯圆盘模空间不一定可定向；相反，从闭黎曼曲面出发的伪全纯映射模空间总是可定向的，而且具有典范定向，这是熟知的事实。此外，由于我们使用伪全纯映射的稳定映射紧化和仓西结构，模空间的基本链一般只能在 $\mathbb Q$ 上定义，而不能在 $\mathbb Z$ 上定义。特别地，我们一般不能改用 $\mathbb Z/2\mathbb Z$ 系数来绕开定向问题。第 8 章将详细研究这些问题。

下面引入“相对 spin 性”，它给出带拉格朗日边界条件的伪全纯圆盘模空间可定向的一个条件。证明见第 8.1 节。

\begin{definition}{3.1.1}\label{def:3-1-1}
\begin{enumerate}[label=\textup{(\arabic*)},leftmargin=2.8em]
  \item 设 $L$ 是辛流形 $M$ 中的定向拉格朗日子流形。如果存在
  $st\in H^2(M;\mathbb Z/2\mathbb Z)$，使得
  \[
    st|_L=w_2(TL)\in H^2(L;\mathbb Z/2\mathbb Z),
  \]
  就称 $L$ 是\emph{相对 spin 拉格朗日子流形}\liftIndex{I026}。

  \item 如果存在 $st\in H^2(M;\mathbb Z/2\mathbb Z)$，同时满足
  \[
    st|_{L^{(0)}}=w_2(TL^{(0)}),
    \qquad
    st|_{L^{(1)}}=w_2(TL^{(1)}),
  \]
  就称 $M$ 中的一对定向拉格朗日子流形
  $(L^{(0)},L^{(1)})$ 是\emph{相对 spin 对}\liftIndex{I027}。
\end{enumerate}
\end{definition}

例如，spin 的 $L$ 显然是相对 spin 的。特别地，当
$\dim_{\mathbb R}M\leq6$ 时，每个可定向拉格朗日子流形都自动是相对 spin 的。要在带拉格朗日边界条件的伪全纯圆盘模空间上指定一个定向，还需要更多数据；我们称之为\emph{相对 spin 结构}。见第 8 章的定义 8.1.2。

除非另有说明，本章始终假定 $L$ 是相对 spin 拉格朗日子流形，并假定
$(L^{(0)},L^{(1)})$ 是相对 spin 对。

本章的第一个主要结果（定理 3.5.11），是为相对 spin 拉格朗日子流形 $L$ 构造一个滤过 $A_\infty$ 代数；它编码所有附着于 $L$ 的伪全纯圆盘的贡献。定义弗洛尔上同调的障碍可以用这个滤过 $A_\infty$ 代数来理解。我们使用万有诺维科夫环
$\Lambda_{0,\mathrm{nov}}(R)$，见第 1 章的约定 (Conv.4)。（在纯代数部分，我们大多对带幺交换环 $R$，例如 $R=\mathbb Z$ 或
$R=\mathbb Z/2\mathbb Z$，使用 $\Lambda_{0,\mathrm{nov}}(R)$；这些代数材料用于 \textbf{[FOOO09I]}。）这个环与文献以及第 2 章中的标准诺维科夫环之间的关系，将在第 5 章第 5.1 节说明。

分别以 $\Lambda_{\mathrm{nov}}^{(k)}$ 和
$\Lambda_{0,\mathrm{nov}}^{(k)}$ 表示它们的 $k$ 次分次部分。需要指明系数环 $R$ 时，写成
$\Lambda_{\mathrm{nov}}(R)$、$\Lambda_{0,\mathrm{nov}}(R)$ 等；有时也写成
$\Lambda_{\mathrm{nov}}^R$、$\Lambda_{0,\mathrm{nov}}^R$。环
$\Lambda_{\mathrm{nov}}$、$\Lambda_{0,\mathrm{nov}}$、
$\Lambda_{0,\mathrm{nov}}^+$ 带有滤过 $F^\lambda\Lambda_{\mathrm{nov}}$ 等，见 (Conv.6) 与 (1.37)。因此，$\Lambda_{\mathrm{nov}}$ 和
$\Lambda_{0,\mathrm{nov}}$ 都成为完备的滤过分次环。注意，
$\Lambda_{\mathrm{nov}}$ 是 $\Lambda_{0,\mathrm{nov}}$ 的一个局部化。

我们将为相对 spin 拉格朗日子流形构造一个
$\Lambda_{0,\mathrm{nov}}(\mathbb Q)$ 上的滤过 $A_\infty$ 代数；在可能时，还会构造取
$\Lambda_{0,\mathrm{nov}}(\mathbb Q)$ 系数的形变弗洛尔上同调。在这一构造中，参数
$\lambda_i$ 在几何上表示伪全纯圆盘的辛面积（或能量），$2n_i$ 表示它的马斯洛夫指标。研究弗洛尔上同调在哈密顿同痕下的不变性时，将使用环
$\Lambda_{\mathrm{nov}}(\mathbb Q)$。

先从经典图景中（无滤过）$A_\infty$ 代数的几何实现开始。在这种情形下，构造相应基本链时可以使用单值截面系，而不必使用多值截面系，所以能够在 $\mathbb Z$ 上构造
$A_\infty$ 代数；细节见第 3.4 节。令 $S^k(L;\mathbb Z)$ 是由 $L$ 的所有余维
$k$ 光滑奇异单形生成的自由 $\mathbb Z$-模，并把
$S(L;\mathbb Z)=\bigoplus_kS^k(L;\mathbb Z)$ 看成上链复形。选取一个由光滑奇异单形组成的可数集 $\mathcal X_L$，并以
$\mathbb Z\mathcal X_L$ 表示 $S(L;\mathbb Z)$ 中由 $\mathcal X_L$ 生成的自由
$\mathbb Z$-子模。第 3.4 节将证明下面的定理（定理 3.4.8）。

取 $L$ 的一个适当光滑单纯剖分，并用一条光滑奇异链 $PD[L]$ 表示基本类。（这里用
$PD[L]$ 而不用 $[L]$，是因为我们把余维 $k$ 的奇异链看成 $k$-上链；这一识别可以视作庞加莱对偶。）

\begin{theorem}{3.1.2}\label{thm:3-1-2}
对任意定向流形 $L$，都存在 $\mathcal X_L$，使得可以在
$\mathbb Z\mathcal X_L$ 上构造一个记作 $\overline{\mathfrak m}^{\,L}$ 的
$\mathbb Z$ 上的 $A_\infty$ 结构。它以 $PD[L]$ 为同伦幺元。此外，
$(\mathbb Z\mathcal X_L,\overline{\mathfrak m}_1^{\,L})$ 的上同调群同 $L$ 的
$\mathbb Z$ 系数上同调群同构。
\end{theorem}

$A_\infty$ 代数同伦幺元的定义将在第 3.3 节给出（定义 3.3.2）。这里必须选取适当的可数集 $\mathcal X_L$，因为我们需要使用横截性论证和贝尔纲定理。

\begin{remark}{3.1.3}\label{rem:3-1-3}
第 4.6 节将证明，上述 $A_\infty$ 代数
$(\mathbb Z\mathcal X_L,\overline{\mathfrak m}^{\,L})$ 的弱同伦等价类由 $L$ 唯一确定。此外，令
$\mathbb R\mathcal X_L=\mathbb Z\mathcal X_L\otimes\mathbb R$，并把
$\overline{\mathfrak m}^{\,L}$ 延拓到 $\mathbb R\mathcal X_L$，则
$(\mathbb R\mathcal X_L,\overline{\mathfrak m}^{\,L})$ 与 $L$ 在 $\mathbb R$ 上的德拉姆 D.G.A. 同伦等价。关于这些结果同 D. Quillen 与丹尼斯·帕内尔·沙利文（Dennis Parnell Sullivan）
\textbf{[Qui69, Sul78]} 所创立的有理同伦论之间的关系，见第 7.5 节的进一步讨论。
\end{remark}

接下来，为拉格朗日子流形 $L$ 构造一个
$\Lambda_{0,\mathrm{nov}}(\mathbb Q)$ 上的滤过 $A_\infty$ 代数。令
$\mathcal X_1(L)$ 是一个由光滑奇异单形组成的可数集，满足
$\mathcal X_1(L)\supset\mathcal X_L$，其中 $\mathcal X_L$ 如定理 3.1.2 中所选。有时把
$\mathcal X_1(L)$ 的元素写成 $P=(|P|,f)$，其中 $|P|$ 是一个
$(n-k)$ 维单形，而 $f\colon|P|\to L$ 是光滑映射。以
$C^k(L;\mathbb Q)$ 表示 $S^k(L;\mathbb Q)$ 中由 $\mathcal X_1(L)$ 生成的子模。定义
$C^\bullet(L;\Lambda_{0,\mathrm{nov}})$ 为
$C^\bullet(L;\mathbb Q)\otimes\Lambda_{0,\mathrm{nov}}$ 的完备化：
\[
  C^\bullet(L;\Lambda_{0,\mathrm{nov}})
  :=\widehat{C^\bullet(L;\mathbb Q)
  \otimes_{\mathbb Q}\Lambda_{0,\mathrm{nov}}},
\]
这里所用的一致结构由类似于 (1.38) 的滤过诱导。对
$x\in C(L;\mathbb Q)$，赋予元素 $xT^\lambda e^n\in
C(L;\Lambda_{0,\mathrm{nov}})$ 次数
\[
  \deg(xT^\lambda e^n):=\deg x+2n,
\]
其中 $\deg x$ 是 $x$ 作为上链的次数。令 $C[1]^k=C^{k+1}$，并以
$\deg'$ 表示移位次数，即
$\deg'(xT^\lambda e^n)=\deg(xT^\lambda e^n)-1$。

令 $\mathcal M_{k+1}^{\mathrm{main}}(\beta)$ 是从亏格零有边稳定曲面到
$(M,L)$ 的伪全纯稳定映射模空间的主分支；其边界上有 $k+1$ 个标记点，同伦类为
$\beta\in\Pi(M;L)=\pi_2(M,L)/\sim$，见定义 2.4.17。对
$P_i=(|P_i|,f_i)\in\mathcal X_1(L)$，考虑纤维积
\begin{equation}
  \mathcal M_{k+1}^{\mathrm{main}}(\beta)
  {}_{(\operatorname{ev}_1,\ldots,\operatorname{ev}_k)}
  \!\times_{f_1\times\cdots\times f_k}
  (P_1\times\cdots\times P_k),
  \tag{3.1.4}\label{eq:3-1-4}
\end{equation}
其中 $\operatorname{ev}_i$（$i=0,1,\ldots,k$）是在第 $i$ 个标记点处的评价映射。这里以及下文，在不会混淆时，用 $P_i$ 代替 $|P_i|$。可以证明~\eqref{eq:3-1-4} 是带仓西结构的空间。把~\eqref{eq:3-1-4} 所示集合记作
$\mathcal M_{k+1}^{\mathrm{main}}(\beta;\vec P)$，其中
$\vec P=(P_1,\ldots,P_k)$。下面是本章的第一个主要定理。

\begin{theorem}{3.1.5}\label{thm:3-1-5}
设 $L$ 是相对 spin 拉格朗日子流形。对任意满足定理 3.1.2 的给定
$\mathcal X_L$，都可以选取 $\mathcal X_1(L)\supset\mathcal X_L$，并在
$\mathcal M_{k+1}^{\mathrm{main}}(\beta;\vec P)$ 上选取一个多值截面系
$\mathfrak s_{\beta,\vec P}$，从而在
$C^\bullet(L;\Lambda_{0,\mathrm{nov}})$ 上构造
$\Lambda_{0,\mathrm{nov}}$ 上的滤过 $A_\infty$ 结构
$\mathfrak m=\{\mathfrak m_k\}$。此外，它有同伦幺元，并满足
$G(L)$-有隙条件\liftIndex{I067}。
\end{theorem}

$G$-有隙条件的定义将在第 3.2 节给出（定义 3.2.26）；它与伪全纯曲线的格罗莫夫紧性有关。设
$G\subset\mathbb R_{\geq0}\times2\mathbb Z$ 是
$\mathbb R_{\geq0}\times2\mathbb Z$ 的一个子幺半群，即
$(0,0)\in G$，并且 $g,g'\in G$ 蕴含 $g+g'\in G$。假定如下条件成立。

\begin{statement}{条件}{3.1.6}\label{cond:3-1-6}
\begin{enumerate}[label=\textup{(3.1.7.\arabic*)},leftmargin=5.6em]
  \item 令 $\pi\colon\mathbb R_{\geq0}\times2\mathbb Z
  \to\mathbb R_{\geq0}$ 为投影，则
  $\pi(G)\subset\mathbb R_{\geq0}$ 是离散的。
  \item $G\cap(\{0\}\times2\mathbb Z)=\{(0,0)\}$。
  \item 对任意 $\lambda$，集合
  $G\cap(\{\lambda\}\times2\mathbb Z)$ 都是有限集。
\end{enumerate}
\end{statement}

在定理 3.1.5 中，$G(L)$ 取如下形式。先定义
\[
  G(L)_0=\bigl\{(\omega[\beta],\mu_L(\beta))\mid
  \beta\in\Pi(L),\ \mathcal M(L;\beta;J)\neq\varnothing\bigr\},
\]
其中 $\omega[\beta]$ 是辛面积，$\mu_L(\beta)$ 是马斯洛夫指标；记号
$\Pi(L)$ 见定义 2.4.17。再定义
\begin{equation}
  G(L)=\mathbb R_{\geq0}\times2\mathbb Z
  \text{ 中由 }G(L)_0\text{ 生成的子幺半群。}
  \tag{3.1.8}\label{eq:3-1-8}
\end{equation}
由格罗莫夫紧性定理\liftIndex{I052}，$G(L)$ 满足条件 3.1.6。（回忆：定向拉格朗日子流形的马斯洛夫指标总是偶数。）一般而言，$G(L)$ 可能依赖于相容殆复结构 $J$。

关于定理中 $\mathcal X_1(L)$ 的选择以及多值截面系
$\mathfrak s_{\beta,\vec P}$ 的构造，还必须解决非平凡的横截性与符号问题。我们把横截性的研究推迟到第 7 章，把符号问题推迟到第 8 章；在此前提下，第 3.5 节将证明定理 3.1.5（定理 3.5.11）。因为构造使用仓西结构和多值截面，各种基本链一般定义在 $\mathbb Q$ 上而非
$\mathbb Z$ 上；这就是此处取 $R=\mathbb Q$ 的原因。

把系数环 $\Lambda_{0,\mathrm{nov}}$ 约化为 $\mathbb Q$，就会像前面所说的那样舍去伪全纯圆盘的贡献。于是得到下述结果，也就是定理 3.5.11 的 (3.5.12.3)；其证明将在第 7.2.8 小节利用第 7.2.9 小节与第 8.9 节的结果完成。

\begin{theorem}{3.1.9}\label{thm:3-1-9}
设 $(C(L;\mathbb Q),\overline{\mathfrak m})$ 是把定理 3.1.5 中的
$(C(L;\Lambda_{0,\mathrm{nov}}),\mathfrak m)$ 约化后得到的
$\mathbb Q$ 上的 $A_\infty$ 代数，并且 $\overline{\mathfrak m}_0=0$。则
$\overline{\mathfrak m}$ 在 $\mathbb Q\mathcal X_L\subset C(L;\mathbb Q)$ 上的限制，等于定理 3.1.2 构造的 $\overline{\mathfrak m}^{\,L}$。
\end{theorem}

因此，滤过 $A_\infty$ 代数
$(C(L;\Lambda_{0,\mathrm{nov}}),\mathfrak m)$ 可以解释为 $L$ 的有理（更准确地说，实）同伦型的量子形变。这里“量子”是指包含来自伪全纯圆盘的所有非线性贡献；物理文献称这些圆盘为\emph{圆盘瞬子}。

由于定理 3.1.5\translatornote{原书此处称“定理 3.1.10 中的算子
$\mathfrak m_1$”，但本节的 (3.1.10) 是紧随其后的方程，并非定理；所指滤过
$A_\infty$ 结构由定理 3.1.5 构造，故改正交叉引用。}所构造的滤过 $A_\infty$ 代数
$(C(L;\Lambda_{0,\mathrm{nov}}),\mathfrak m)$ 中的
$\mathfrak m_0$ 可能非零，算子 $\mathfrak m_1$ 一般不满足
$\mathfrak m_1\circ\mathfrak m_1=0$；这正是定义 $L$ 的弗洛尔上同调所遇到的障碍。另一方面，利用任意
$b\in C[1]^0=C^1$，可把 $\mathfrak m$ 形变为
$\mathfrak m^b$：
\begin{align*}
  \mathfrak m_k^b(x_1,\ldots,x_k)
  &={}
  \sum_{\ell_0,\ldots,\ell_k\geq0}
  \mathfrak m_{k+\sum_i\ell_i}
  \bigl(\\[-0.35em]
  &\qquad
    \underbrace{b,\ldots,b}_{\ell_0},x_1,
    \underbrace{b,\ldots,b}_{\ell_1},\ldots,\\[-0.15em]
  &\qquad
    \underbrace{b,\ldots,b}_{\ell_{k-1}},x_k,
    \underbrace{b,\ldots,b}_{\ell_k}
  \bigr)\\
  &={}
  \mathfrak m(e^b x_1e^b\cdots e^b x_ke^b),
\end{align*}
其中 $k=0,1,2,\ldots$；于是 $(C,\mathfrak m^b)$ 仍是滤过
$A_\infty$ 代数，见定义 3.6.9。这里采用记号
\[
  e^b:=1+b+b\otimes b+b\otimes b\otimes b+\cdots.
\]
进一步，若 $b$ 满足 Maurer--Cartan 方程的 $A_\infty$ 版本
\liftIndex{I068}\liftIndex{I078}
\begin{equation}
  \mathfrak m(e^b):=
  \mathfrak m_0(1)+\mathfrak m_1(b)+\mathfrak m_2(b,b)
  +\mathfrak m_3(b,b,b)+\cdots=0,
  \tag{3.1.10}\label{eq:3-1-10}
\end{equation}
就有 $\mathfrak m_0^b=0$（命题 3.6.10）。这时
$\mathfrak m_1^b\circ\mathfrak m_1^b=0$，因而可以定义
$(C(L;\Lambda_{0,\mathrm{nov}}),\mathfrak m_1^b)$ 的上同调。称
Maurer--Cartan 方程~\eqref{eq:3-1-10} 的解 $b$ 为\emph{定界余链}
\liftIndex{I004}\liftIndex{I009}。（为使~\eqref{eq:3-1-10} 左端收敛，还假定
$b\equiv0\pmod{\Lambda_{0,\mathrm{nov}}^+}$。）

把解集记作
\[
  \widehat{\mathcal M}
  \bigl((C(L;\Lambda_{0,\mathrm{nov}}),\mathfrak m)\bigr)
  =\widehat{\mathcal M}(L).
\]
若 $\widehat{\mathcal M}(L)\neq\varnothing$，就称 $L$ 是\emph{无障碍的}
\liftIndex{I069}\liftIndex{I110}\liftIndexRangeStart{I075}。若 $L$ 无障碍，则对
$b\in\widehat{\mathcal M}(L)$ 定义由 $b$ 形变的弗洛尔上同调
\[
  HF(L,b;\Lambda_{0,\mathrm{nov}})
  :=H(C(L;\Lambda_{0,\mathrm{nov}}),\mathfrak m_1^b).
\]

下一个问题是寻找保证 $L$ 无障碍的条件。为此，我们将构造
$H(L;\mathbb Q)$ 中的一列元素
$\{[o_k(L)]\}_{k=1,2,\ldots}$，称为\emph{（弗洛尔）障碍类}
\liftIndexRangeStart{I073}\liftIndexRangeStart{I348}。由于 $G(L)$ 满足条件 3.1.6，可以写成
\[
  G(L)=\bigl\{(\lambda_i,\mu_i)
  =(\omega[\beta_i],\mu_L(\beta_i))\mid i=0,1,2,\ldots\bigr\},
\]
并使 $i<j$ 蕴含 $\lambda_i\leq\lambda_j$。（注意，
$\lambda_0=\mu_0=0$ 且 $0<\lambda_1$。）若某些
$\beta,\beta'$ 满足 $\omega[\beta]=\omega[\beta']$，则 $G(L)$ 的这种枚举不唯一；不过所得结果与枚举无关。

\begin{theorem}{3.1.11}\label{thm:3-1-11}
设 $(C(L;\Lambda_{0,\mathrm{nov}}),\mathfrak m)$ 是定理 3.1.5 构造的滤过
$A_\infty$ 代数。则存在上链序列 $o_k(L)$ 和 $b_k(L)$（$k=1,2,\ldots$），具有下列性质。

\smallskip
\noindent\textup{(3.1.12.1)}
$[o_k(L)]\in H^{2-\mu_k}(L;\mathbb Q)$，并且
$b_k(L)\in C^{1-\mu_k}(L;\mathbb Q)$。这里
$\mu_k=\mu_L(\beta_k)$ 是 $\beta_k$ 的马斯洛夫指标，而且是偶数。

\noindent\textup{(3.1.12.2)}
若对所有满足 $\lambda_j<\lambda_k$ 的 $j$，$b_j(L)$ 和
$o_j(L)$ 均已定义，则闭上链 $o_k(L)$ 有定义；
$[o_k(L)]$ 依赖于这些 $b_j(L)$ 和 $o_j(L)$。

\noindent\textup{(3.1.12.3)}
若 $[o_k(L)]=0$，则把 $b_k(L)$ 定义为满足
$-\mathfrak m_{1,0}(b_k)=o_k$ 的一个上链。

\noindent\textup{(3.1.12.4)}
$L$ 无障碍，当且仅当可以归纳地选取 $b_k(L)$，使所有
$[o_k(L)]$ 都为零。
\end{theorem}

在任意满足 $G$-有隙条件的抽象滤过 $A_\infty$ 代数中，都能纯代数地构造这些障碍类。事实上，第 3.6.2 小节将在这一一般框架下完成构造（定理 3.6.18）。

还可以定义障碍类的一个细化，记作 $[o_k(L;\mathrm{weak})]$；此时只假定马斯洛夫指标非正的先行障碍类消失，见定理 3.6.43。这些细化障碍类同下述定义弗洛尔上同调的较弱条件有关。

回忆：我们的 $A_\infty$ 代数
$(C(L;\Lambda_{0,\mathrm{nov}}),\mathfrak m)$ 有同伦幺元。第 3.6.3 小节将对具有同伦幺元的 $A_\infty$ 代数引入一个较弱的条件，用以定义弗洛尔上同调。设
$(C,\mathfrak m)$ 是 $\Lambda_{0,\mathrm{nov}}$ 上以 $\mathbf e$ 为幺元的
$A_\infty$ 代数，见定义 3.2.20(3)。对 $b\in C[1]^0$，考虑比
Maurer--Cartan 方程更弱的方程
\begin{equation}
  \mathfrak m(e^b)=ce\mathbf e,
  \tag{3.1.13}\label{eq:3-1-13}
\end{equation}
其中某个 $c\in\Lambda_{0,\mathrm{nov}}^{+(0)}$。（这里
$\deg\mathbf e=0$ 而 $\deg e=2$，所以
$\deg(ce\mathbf e)-1=+1$。）令
\[
  \widehat{\mathcal M}_{\mathrm{weak}}(C)
  =\bigl\{b\in C[1]^0\mid
  \mathfrak m(e^b)=ce\mathbf e
  \text{，其中某个 }c\in\Lambda_{0,\mathrm{nov}}^{+(0)}\bigr\}.
\]
如果 $\widehat{\mathcal M}_{\mathrm{weak}}(C)\neq\varnothing$，就称以
$\mathbf e$ 为幺元的 $\Lambda_{0,\mathrm{nov}}$ 上 $A_\infty$ 代数\emph{弱无障碍}
\liftIndex{I070}\liftIndex{I353}。对具有同伦幺元的 $A_\infty$ 代数
$(C,\mathfrak m)$，有一个以 $\mathbf e^+$ 为幺元的相伴
$A_\infty$ 代数 $(C^+,\mathfrak m^+)$（定义 3.3.2）。若
$\widehat{\mathcal M}_{\mathrm{weak}}(C^+)\neq\varnothing$，就称具有同伦幺元的
$(C,\mathfrak m)$ 弱无障碍\liftIndex{I076}\liftIndexRangeEnd{I075}。

对 $b\in\widehat{\mathcal M}_{\mathrm{weak}}(C^+)$，定义
\[
  \delta_{b,b}(x)
  =\sum_{k_1,k_2\geq0}
  \mathfrak m_{k_1+k_2+1}^+
  (\underbrace{b,\ldots,b}_{k_1},x,
   \underbrace{b,\ldots,b}_{k_2}),
\]
这正是 $\mathfrak m_1^{+b}(x)$。由引理 3.6.33，对
$b\in\widehat{\mathcal M}_{\mathrm{weak}}(C^+)$ 有
$\delta_{b,b}\circ\delta_{b,b}=0$。因此，若
$(C(L;\Lambda_{0,\mathrm{nov}}),\mathfrak m)$ 弱无障碍，仍可对
\[
  b\in\widehat{\mathcal M}_{\mathrm{weak}}(L)
  :=\widehat{\mathcal M}_{\mathrm{weak}}
  (C(L;\Lambda_{0,\mathrm{nov}})^+)
\]
定义由 $b$ 形变的弗洛尔上同调
\[
  HF(L,b;\Lambda_{0,\mathrm{nov}})
  :=H(C(L;\Lambda_{0,\mathrm{nov}})^+,\mathfrak m_1^{+b}).
\]
定理 3.6.43 表明，如果所有细化障碍类
$[o_k(L;\mathrm{weak})]$ 都消失，那么 $L$ 弱无障碍。

现在还允许伪全纯圆盘除边界标记点外带有一个内部标记点；第 3.8 节将由此证明下面的结果，见定理 3.8.11。\liftIndexRangeStart{I074}

\begin{theorem}{3.1.14}\label{thm:3-1-14}
设 $L\subset M$ 是维数为 $n$ 的相对 spin 拉格朗日子流形，则：
\begin{enumerate}[label=\textup{(\arabic*)},leftmargin=2.8em]
  \item 若 $[GW_{0,1}(M)(L)]\neq0$，则
  $\widehat{\mathcal M}(L)=\varnothing$；
  \item 若 $[GW_{0,1}(M)(L)]=0$，则障碍类 $[o_k(L)]$ 属于
  Gysin 同态
  \[
    i_!\colon H^{2-\mu_k}(L;\mathbb Q)
    \longrightarrow H^{n+2-\mu_k}(M;\mathbb Q)
  \]
  的核。
\end{enumerate}
\end{theorem}

这里
\[
  [GW_{0,1}(M)(L)]
  \in H(M;\mathbb Q)\otimes\Lambda_{0,\mathrm{nov}}
\]
是带两个标记点的亏格零格罗莫夫--威腾不变量，其中一个标记点附着于
$L$。精确定义见第 3.8.1 小节；若干应用见第 3.8.2 小节的推论
3.8.16--3.8.19。\liftIndexRangeEnd{I074}

\begin{theorem}{3.1.15}\label{thm:3-1-15}
存在一列障碍类
\[
  [o_k(L;\mathrm{def})]
  \in
  \frac{H^{2-\mu_k}(L;\mathbb Q)}
  {\operatorname{Im}\bigl(i^*\colon
    H^{2-\mu_k}(M;\mathbb Q)
    \to H^{2-\mu_k}(L;\mathbb Q)\bigr)},
\]
以及
\[
  b_k\in C^{1-\mu_k}(L;\mathbb Q),
  \qquad
  \mathfrak b_k\in C^{2-\mu_k}(M;\mathbb Q),
  \qquad
  \delta_M\mathfrak b_k=0,
\]
它们具有下列性质。\translatornote{原书在定理 3.1.15 中把第一项印成
$b_k\in C^{1-\mu_k}_2(L;\mathbb Q)$。定理 3.8.41 的正式重述写作
$b_k\in C^{1-\mu_k}(L;\mathbb Q)$，尽管其证明稍后又重复出现了下标 $2$；本章并未定义与此处相容的
$C_2(L)$。故依正式定理陈述与类型关系删去下标 $2$。}

\smallskip
\noindent\textup{(3.1.16.1)}
若对所有满足 $\lambda_j<\lambda_k$ 的 $j$，$b_j$、
$\mathfrak b_j$ 和 $o_j$ 均已定义，则闭上链
$o_k(L;\mathrm{def})$ 有定义；$[o_k(L;\mathrm{def})]$ 依赖于这些
$b_j$、$\mathfrak b_j$ 和 $o_j$。

\noindent\textup{(3.1.16.2)}
若 $[o_k(L;\mathrm{def})]=0$，则 $b_k$、$\mathfrak b_k$ 有定义。

\noindent\textup{(3.1.16.3)}
如果所有障碍类 $[o_k(L;\mathrm{def})]$ 都有定义且均为零，那么
$L$ 经 bulk 形变后无障碍
\liftIndex{I071}\liftIndex{I072}\liftIndex{I077}\liftIndex{I354}。
\end{theorem}

\begin{corollary}{3.1.17}\label{cor:3-1-17}
如果对每个 $k$，
\[
  i^*\colon H^{2k}(M;\mathbb Q)\longrightarrow H^{2k}(L;\mathbb Q)
\]
都是满射，那么 $L$ 经 bulk 形变后无障碍。
\end{corollary}

如果 $L$ 经 bulk 形变后无障碍，还可以定义由
$(b,\mathfrak b)\in\widehat{\mathcal M}_{\mathrm{def}}(L)$ 形变的弗洛尔上同调（引理 3.8.60）；集合
$\widehat{\mathcal M}_{\mathrm{def}}(L)$ 见定义 3.8.40。把幺元或同伦幺元纳入后，同样可以定义
$\widehat{\mathcal M}_{\mathrm{weak,def}}(L)$ 以及一列障碍闭上链
$o_k(L;\mathrm{weak,def})$，由此推广定理 3.1.15；见第 3.8.5 小节。
\liftIndexRangeEnd{I073}\liftIndexRangeEnd{I348}

至此讨论的是单个拉格朗日子流形。下面考虑一对相对 spin 拉格朗日子流形
$(L^{(0)},L^{(1)})$。

在第 3.7.4 小节，先假定 $L^{(0)}$ 与 $L^{(1)}$ 横截。对这样的对子，定义一个分次滤过
$\Lambda_{0,\mathrm{nov}}$-模
$C(L^{(1)},L^{(0)};\Lambda_{0,\mathrm{nov}})$（见定义 3.7.20）。另一方面，对每个
$L^{(i)}$（$i=0,1$），已有滤过 $A_\infty$ 代数
$(C(L^{(i)};\Lambda_{0,\mathrm{nov}}),\mathfrak m^{(i)})$。我们将在
$C(L^{(1)},L^{(0)};\Lambda_{0,\mathrm{nov}})$ 上构造滤过 $A_\infty$ 双模结构（定理 3.7.21）。

\begin{theorem}{3.1.18}\label{thm:3-1-18}
设 $(L^{(0)},L^{(1)})$ 是一对彼此横截的相对 spin 拉格朗日子流形。则
$C(L^{(1)},L^{(0)};\Lambda_{0,\mathrm{nov}})$ 上存在一个左
$(C(L^{(1)};\Lambda_{0,\mathrm{nov}}),\mathfrak m^{(1)})$、右
$(C(L^{(0)};\Lambda_{0,\mathrm{nov}}),\mathfrak m^{(0)})$ 的滤过
$A_\infty$ 双模结构。

它是 $G(L^{(1)},L^{(0)})$-有隙的。此外，
上述左右两个代数的一对同伦幺元
$\{\mathbf e_1,\mathbf e_0\}$ 作为该双模的同伦幺元起作用。
\end{theorem}

定理中各种概念的精确定义将在第 3.7.1 小节给出；
$G(L^{(1)},L^{(0)})$ 的定义另见 (3.7.43)。还可以证明下面的结果（定理 3.7.45）。

\begin{theorem}{3.1.19}\label{thm:3-1-19}
设 $(L^{(0)},L^{(1)})$ 是一对彼此横截的相对 spin 拉格朗日子流形。则对每个
$b_0\in\widehat{\mathcal M}(L^{(0)})$、
$b_1\in\widehat{\mathcal M}(L^{(1)})$，都存在上同调
\[
  HF((L^{(1)},b_1),(L^{(0)},b_0);\Lambda_{0,\mathrm{nov}}),
\]
称为由 $(b_0,b_1)$ 形变的 $(L^{(0)},L^{(1)})$ 的弗洛尔上同调。
\end{theorem}

令 $\Omega(L^{(0)},L^{(1)})$ 是连接 $L^{(0)}$ 与 $L^{(1)}$ 的所有道路
$\ell\colon[0,1]\to M$ 之空间。若 $L^{(0)}$、$L^{(1)}$ 连通，则
\[
  \pi_0(\Omega(L^{(0)},L^{(1)}))
  =i_*(\pi_1(L^{(0)}))\backslash\pi_1(M)
  /i_*(\pi_1(L^{(1)})).
\]
每个点 $p\in L^{(0)}\cap L^{(1)}$ 都确定一条常值道路
$t\mapsto p$，从而确定 $\pi_0(\Omega(L^{(0)},L^{(1)}))$ 的一个元素。由此得到弗洛尔上同调群的分解（见注 3.7.46）：
\begin{equation}
\begin{split}
  &HF((L^{(1)},b_1),(L^{(0)},b_0);\Lambda_{0,\mathrm{nov}})\\
  &\qquad\cong
  \bigoplus_{[\ell_0]\in\pi_0(\Omega(L^{(0)},L^{(1)}))}
  HF((L^{(1)},b_1),(L^{(0)},b_0);\ell_0;\Lambda_{0,\mathrm{nov}}).
\end{split}
  \tag{3.1.20}\label{eq:3-1-20}
\end{equation}

要注意，若 $L^{(0)}$ 与 $L^{(1)}$ 弱无障碍，一般仍不能对
$b_i\in\widehat{\mathcal M}_{\mathrm{weak}}(L^{(i)})$ 定义由 $b_i$ 形变的对子
$(L^{(0)},L^{(1)})$ 的弗洛尔上同调；但对每个 $i=0,1$，单个
$L^{(i)}$ 的弗洛尔上同调
$HF(L^{(i)},b_i;\Lambda_{0,\mathrm{nov}})$ 都有定义。这一点对许多应用都很重要。为了描述对子形变弗洛尔上同调有定义的条件，引入称为\emph{势函数}
\liftIndex{I079}的映射
\[
  \mathfrak{PO}\colon
  \widehat{\mathcal M}_{\mathrm{weak}}(L)
  \longrightarrow\Lambda_{0,\mathrm{nov}}^{+(0)},
\]
其定义是：若~\eqref{eq:3-1-13} 中的常数为 $c$，则
$\mathfrak{PO}(b)=c$。于是可证明如下结果（命题 3.7.17）。

\begin{theorem}{3.1.21}\label{thm:3-1-21}
设 $(L^{(0)},L^{(1)})$ 是一对彼此横截的相对 spin 拉格朗日子流形。对每个
\[
  b_0\in\widehat{\mathcal M}_{\mathrm{weak}}(L^{(0)}),
  \qquad
  b_1\in\widehat{\mathcal M}_{\mathrm{weak}}(L^{(1)})
\]
满足 $\mathfrak{PO}(b_0)=\mathfrak{PO}(b_1)$，都能定义由
$b_0,b_1$ 形变的弗洛尔上同调
\[
  HF((L^{(1)},b_1),(L^{(0)},b_0);\Lambda_{0,\mathrm{nov}}).
\]
\end{theorem}

它同样具有类似于~\eqref{eq:3-1-20} 的分解。还可以用同样的方式定义
\[
  \mathfrak{PO}\colon
  \widehat{\mathcal M}_{\mathrm{weak,def}}(L)
  \longrightarrow\Lambda_{0,\mathrm{nov}}^{+(0)}.
\]

第 3.7.5 小节将研究 $L^{(0)}$ 与 $L^{(1)}$ 干净相交的情形，即弗洛尔上同调的博特--莫尔斯版本。交集
$L^{(1)}\cap L^{(0)}$ 的各个连通分支都是光滑流形，但不一定可定向。仔细处理定向层之后，可以得到与横截情形相同的结论。

\begin{theorem}{3.1.22}\label{thm:3-1-22}
设 $(L^{(0)},L^{(1)})$ 是相对 spin 对，并假定 $L^{(0)}$ 与
$L^{(1)}$ 干净相交。则定理 3.1.18、3.1.19 和 3.1.21 的结论仍然成立。
\end{theorem}

滤过 $A_\infty$ 代数
$(C(L;\Lambda_{0,\mathrm{nov}}),\mathfrak m)$ 本身可以看成一个以自身为左、右代数的滤过
$A_\infty$ 双模，见例 3.7.6。另一方面，在定理 3.1.22 中令
$L^{(0)}=L^{(1)}=:L$，由干净相交构造也得到一个以该代数为左、右代数的滤过
$A_\infty$ 双模。若构造双模结构时使用的殆复结构族
$\{J_t\}_t$ 与 $t$ 无关，则这两个 $A_\infty$ 双模可以识别（命题 3.7.75）。所以，当 $L$ 无障碍或弱无障碍时，
\[
  HF(L,b;\Lambda_{0,\mathrm{nov}})
  \cong HF((L,b),(L,b);\Lambda_{0,\mathrm{nov}}).
\]

进一步，在定理 3.1.22 中考虑
$L^{(1)}=\psi(L^{(0)})$、其中 $\psi$ 是哈密顿微分同胚的情形，就得到形变弗洛尔上同调
\[
  HF((\psi(L),\psi_*(b)),(L,b);\Lambda_{0,\mathrm{nov}}).
\]
$\psi_*(b)$ 见第 4.1 节定理 4.1.3。第 3.7.6 小节的例子将表明，这个上同调一般依赖于 $\psi$。不过，第 5 章将证明
\[
\begin{split}
  &HF((\psi(L),\psi_*(b)),(L,b);\Lambda_{\mathrm{nov}})\\
  &\quad:=HF((\psi(L),\psi_*(b)),(L,b);\Lambda_{0,\mathrm{nov}})
  \otimes_{\Lambda_{0,\mathrm{nov}}}\Lambda_{\mathrm{nov}}
\end{split}
\]
与哈密顿微分同胚 $\psi$ 无关。合并这些结论，得到
\[
  HF(L,b;\Lambda_{\mathrm{nov}})
  \cong HF((\psi(L),\psi_*(b)),(L,b);\Lambda_{\mathrm{nov}}).
\]

本章安排如下。第 3.2 节引入滤过 $A_\infty$ 代数与 $A_\infty$ 同态；第 3.3 节引入同伦幺元，以及含幺或同伦含幺的 $A_\infty$ 同态。这两节完全在代数框架中展开，不出现拉格朗日子流形。随后，第 3.4 节在每个定向流形 $L$ 上几何实现带同伦幺元的经典
$A_\infty$ 代数 $(\mathbb Z\mathcal X_L,\overline{\mathfrak m}^{\,L})$。第 3.5 节为相对 spin 拉格朗日子流形 $L$ 构造带同伦幺元的滤过（量子）
$A_\infty$ 代数
$(C(L;\Lambda_{0,\mathrm{nov}}),\mathfrak m)$。要强调的是，第 3.4、3.5 节始终不作任何无障碍性假设。

第 3.6 节在滤过 $A_\infty$ 代数的框架下引入定界余链和
Maurer--Cartan 方程，并构造障碍类；还定义无障碍与弱无障碍，介绍势函数
$\mathfrak{PO}$，并研究弱无障碍条件同 $L$ 的弗洛尔上同调存在性之间的关系。

第 3.7 节考虑一对相对 spin 拉格朗日子流形
$(L^{(0)},L^{(1)})$，利用前面各节的构造形变
$(L^{(0)},L^{(1)})$ 的标准弗洛尔上边界算子。为此，第 3.7.1、3.7.2 小节分别在分次滤过
$\Lambda_{0,\mathrm{nov}}$-模上引入滤过 $A_\infty$ 双模结构和
$A_\infty$ 双模同态的代数概念，第 3.7.3 小节讨论弱无障碍性。第 3.7 节的主要任务，是在
$C(L^{(1)},L^{(0)};\Lambda_{0,\mathrm{nov}})$ 上构造左
$(C(L^{(1)};\Lambda_{0,\mathrm{nov}}),\mathfrak m^{(1)})$、右
$(C(L^{(0)};\Lambda_{0,\mathrm{nov}}),\mathfrak m^{(0)})$ 的滤过
$A_\infty$ 双模结构。第 3.7.4 小节处理 $L^{(0)}$ 与 $L^{(1)}$ 横截的情形，第 3.7.5 小节处理二者干净相交的情形。

定义这个 $A_\infty$ 双模结构时，不必假定
上述左右两个滤过 $A_\infty$ 代数无障碍或弱无障碍；换言之，任意相对 spin 对都有相应的滤过
$A_\infty$ 双模。若这两个代数无障碍，则可分别为
$L^{(0)}$、$L^{(1)}$ 选取定界余链 $b_0,b_1$，再用 $A_\infty$ 双模结构得到形变弗洛尔上链复形
\[
  (C(L^{(1)},L^{(0)};\Lambda_{0,\mathrm{nov}}),\delta_{b_1,b_0}).
\]
第 3.7.6 小节将给出若干简单的弗洛尔上同调计算例子。

第 3.8 节还考虑带内部标记点的伪全纯圆盘模空间。利用从有边黎曼曲面出发、带一个内部标记点而无边界标记点的伪全纯映射模空间，我们将构造一列算子
$\mathfrak p_k$ 并研究其性质。借助算子 $\mathfrak p$，研究第 3.6.2、3.6.3 小节构造的障碍类，并在第 3.8.1--3.8.3 小节给出若干保证（弱）无障碍的有用条件。

进一步，第 3.8.4 小节考虑带若干内部标记点的情形，构造一列算子
$\mathfrak q$。这些算子用 $H^2(M;\Lambda_{0,\mathrm{nov}}^+)$ 的元素给出滤过
$A_\infty$ 代数的 bulk 形变。第 3.8.5 小节还将构造障碍类的相对版本。算子
$\mathfrak q$ 的构造梗概见第 3.8.6 小节。这一思想也用于第 3.8.7 小节的滤过
$A_\infty$ 双模，并给出一列算子 $\mathfrak r$；其构造在第 3.8.8 小节说明。这些算子的构造细节见第 7.4 节。

最后一个小节把 $\mathfrak p_k$ 的构造同 $\mathfrak q$ 的构造结合起来，得到一列推广算子
$\mathfrak p_{\ell,k}$，并讨论其若干性质。第 6 章的谱序列论证将使用这些算子。第 3.8 节还将证明若干公式，给出算子
$\mathfrak p,\mathfrak q,\mathfrak r$ 的基本性质；第 7.4 节会用
$L_\infty$ 结构的语言重新解释这些公式，并给出完整构造。

为了从几何上构造 $A_\infty$ 代数，必须解决若干横截性问题；第 7 章将更详细地讨论这一点。符号（或定向）也需要仔细处理；第 8 章将作出精确论证。因此，本章只严格陈述结果，并在暂不处理横截性和符号问题的前提下给出证明。

还要指出，本章构造需要选定并固定各种结构，例如相容殆复结构、仓西结构、多值截面、三角剖分以及可数集
$\mathcal X_L$、$\mathcal X_1(L)$ 等。必须研究构造对这些选择的依赖或独立性；这将在后续第 4、5 章完成。
\translatornote{原书此处把第二个可数集印作未定义的
$\mathcal X(L)$；本节此前及第 3.5 节反复定义并使用的是
$\mathcal X_1(L)$，故补回下标 $1$。}
